\documentclass[12pt]{article}
\usepackage[utf8]{inputenc}
\usepackage[T1]{fontenc}
\usepackage{hyperref}
\title{San Mart\'in Itunyoso Triqui Language Resources for K-12 Curriculum Development}
\author{Compiled by Claude Sonnet 5 (claude-sonnet-5) for Andr\'as Kornai}
\date{}
\begin{document}
\maketitle

\section{Introduction}
San Mart\'in Itunyoso Triqui (ISO 639-3 trq) is an endangered Otomanguean language of Oaxaca, Mexico. It is tonal. Most of the material found comes from the work of Christian DiCanio.

\section{Primary Educational Resources}
\begin{itemize}
\item \emph{Libro pedag\'ogico Triqui} (2019), a reading booklet for the San Mart\'in Itunyoso alphabet, co-created by DiCanio with Basileo and Wilberto Mart\'inez Cruz: \url{https://www.acsu.buffalo.edu/~cdicanio/pdfs/Libro_pedagogico_Triqui_2019.pdf}
\end{itemize}

\section{Linguistic and Reference Materials}
\begin{itemize}
\item DiCanio's publications page, which lists the Triqui--Spanish dictionary (online version), an in-progress reference grammar, and phonology papers: \url{https://www.acsu.buffalo.edu/~cdicanio/publications.html}
\item DiCanio (2008), \emph{The Phonetics and Phonology of San Mart\'in Itunyoso Trique}, Berkeley dissertation: \url{https://glottolog.org/resource/reference/id/159284}
\item DiCanio (2010), Illustrations of the IPA: San Mart\'in Itunyoso Trique: \url{https://glottolog.org/resource/reference/id/548421}
\item DiCanio (2016), chapter on tonal classes in person morphology: \url{https://glottolog.org/resource/reference/id/548423}
\end{itemize}
No finished published grammar was found.

\section{Religious and Community Resources}
None were found in these searches.

\section{Digital and Audio Resources}
\begin{itemize}
\item AILLA (Archive of the Indigenous Languages of Latin America) collection, Itunyoso Triqui: \url{https://islandora-ailla.lib.utexas.edu/islandora/object/ailla%253A244044}
\item AILLA: lexical elicitation recordings (2013), including the first 460 words of the dictionary: \url{https://islandora-ailla.lib.utexas.edu/islandora/object/ailla%3A244017}
\end{itemize}

\section{Curriculum Development Context}
The reference description says the language is not known to be taught in schools. The 2019 primer and the dictionary are the main starting points for classroom materials. Beyond them, little was found, and any curriculum would need to be developed with the community.
\end{document}